%%%%%%%%%%%%%%%%%
% CV tailored for Scalian - Ingénieur Validation Systèmes Embarqués Défense
%%%%%%%%%%%%%%%%%

\documentclass[8pt,a4paper,ragged2e,withhyper]{../_shared/altacv}

\geometry{left=1.25cm,right=1.25cm,top=.5cm,bottom=1.cm,columnsep=1.2cm}

\usepackage{paracol}
\usepackage{fontawesome5}
\usepackage{hyperref}

\iftutex 
  \setmainfont{Roboto Slab}
  \setsansfont{Lato}
  \renewcommand{\familydefault}{\sfdefault}
\else
  \usepackage[rm]{roboto}
  \usepackage[defaultsans]{lato}
  \renewcommand{\familydefault}{\sfdefault}
\fi

\definecolor{SlateGrey}{HTML}{2E2E2E}
\definecolor{LightGrey}{HTML}{666666}
\definecolor{DarkPastelRed}{HTML}{8AC0D9}
\definecolor{PastelRed}{HTML}{00B0DA}
\definecolor{GoldenEarth}{HTML}{B4AD0C}

\colorlet{name}{black}
\colorlet{tagline}{PastelRed}
\colorlet{heading}{DarkPastelRed}
\colorlet{headingrule}{GoldenEarth}
\colorlet{subheading}{PastelRed}
\colorlet{accent}{PastelRed}
\colorlet{emphasis}{SlateGrey}
\colorlet{body}{LightGrey}

\definecolor{violet}{HTML}{5A2582}
\definecolor{rouge}{HTML}{F53B3B}
\definecolor{noir}{HTML}{00232C}
\definecolor{bleu}{HTML}{00B0DA}
\definecolor{rose}{HTML}{F831A0}
\definecolor{jaune}{HTML}{FFEC00}
\definecolor{vert}{HTML}{34AD6C}

\renewcommand{\namefont}{\Huge\rmfamily\bfseries}
\renewcommand{\personalinfofont}{\footnotesize}
\renewcommand{\cvsectionfont}{\LARGE\rmfamily\bfseries}
\renewcommand{\cvsubsectionfont}{\large\bfseries}
\renewcommand{\cvItemMarker}{{\small\textbullet}}
\renewcommand{\cvRatingMarker}{\faCircle}

\input{../_shared/pubs-num.tex}
\addbibresource{../_shared/sample.bib}

\begin{document}

\name{BOILEAU Guillaume}
\tagline{Ingénieur IVV / Validation Systèmes | Python, Automatisation de Tests, Analyse d'Anomalies \& Systèmes Complexes}

\photoL{2.8cm}{../_shared/moi}

\personalinfo{%
  \email{guillaume.boileaupro@gmail.com}
  \phone{+33 6 59 94 85 84}
  \location{Alpes-Maritimes, FRANCE}
  \linkedin{guillaume-boileau-phd-891b81265}
  \researchgate{Guillaume-Boileau-2}
  \orcid{0000-0002-3576-6968}
}

\makecvheader
\vspace{-0.3cm}

\AtBeginEnvironment{itemize}{\small}
\columnratio{0.64}

\begin{paracol}{2}

\cvsection{Experience}

\cvevent{\textbf{Software Engineer - System Integration, Production Diagnostics \& Hardware Interfaces}}{VEV Platform Services France}{2025 -- 2026}{Valbonne, France}

\textbf{Context:} Development and integration of software services for EV charging infrastructure within a CPMS platform connected to operational hardware systems.

\textbf{Objective:} Support robust software integration, operational data flows and interface reliability between platform services and charging station systems.

\begin{itemize}
  \item \textbf{Software development:} Developed web and backend services using TypeScript, Angular and Node.js in an operational industrial environment.
  \item \textbf{System integration:} Contributed to the integration and maintenance of software services interacting with field hardware and charging infrastructure.
  \item \textbf{Hardware/software interfaces:} Analysed interactions between platform services, operational data flows and charging station systems.
  \item \textbf{Operational validation:} Checked data consistency, service behaviour and interface continuity in production-like conditions.
  \item \textbf{Incident analysis:} Investigated production issues involving software behaviour, data exchanges and hardware-connected systems.
  \item \textbf{First-level diagnostics:} Analysed abnormal behaviours, communication issues and data inconsistencies before escalation or correction follow-up.
  \item \textbf{Data exchange:} Implemented and monitored FTP-based operational data exchange services.
  \item \textbf{Agile tooling:} Used Zenhub to support task tracking, backlog follow-up, iterative development and technical coordination.
\end{itemize}

\divider

\cvevent{\textbf{Senior R\&D Engineer - Space Systems, IVV, Test Automation \& Validation Workflows}}{Laboratoire Lagrange, Observatoire de la Côte d'Azur, CNES}{2023 -- 2025}{Nice, France}

\textbf{Context:} Engineering activities for the LISA ESA/NASA space mission, involving complex software chains, model integration, simulation workflows, validation campaigns and international coordination.

\textbf{Objective:} Develop, integrate and validate software and simulation workflows under strong consistency, traceability, performance and verification constraints.

\begin{itemize}
  \item \textbf{Space systems:} Contributed to mission-level analysis workflows for the LISA ESA/NASA space mission in a large international technical collaboration.
  \item \textbf{V-cycle logic:} Supported engineering activities aligned with specification, integration, verification and validation logic for complex mission-level software workflows.
  \item \textbf{IVV logic:} Structured integration, verification and validation workflows for complex models, software chains and simulation outputs.
  \item \textbf{Python automation:} Developed Python tools to automate model comparison, consistency checks, validation workflows and technical analysis tasks.
  \item \textbf{Test strategy:} Defined comparison campaigns, validation checks and acceptance logic for technical outputs produced by multiple contributors.
  \item \textbf{Non-regression logic:} Compared successive outputs, configurations and model behaviours to detect deviations, regressions and unexpected changes.
  \item \textbf{Script execution:} Executed analysis and validation scripts, reviewed outputs and consolidated results for technical interpretation.
  \item \textbf{Anomaly analysis:} Investigated inconsistencies, deviations and unexpected behaviours affecting result reliability and interpretation.
  \item \textbf{Software platform:} Developed CATpyLISA, a Python platform for large-scale model integration, comparison, validation and technical review.
  \textcolor{DarkPastelRed}{\href{https://gitlab.oca.eu/gboileau/lisa_l3_tools}{\bf GitLab:CATpyLISA}}
  \item \textbf{Technical reporting:} Produced validation notes, technical documentation, analysis summaries and evidence supporting technical reviews.
  \item \textbf{Technical coordination:} Coordinated contributors, supervised interns and supported validation activities across multiple stakeholders.
  \item \textbf{Documentation tooling:} Used Confluence to support technical documentation, coordination notes, validation follow-up and knowledge sharing within project activities.
\end{itemize}

\divider

\cvevent{\textbf{Data Scientist - Diagnostics, Test Analysis \& Performance Validation}}{Universiteit Antwerpen, Belgium}{2021 -- 2023}{Antwerp, Belgium}

\textbf{Context:} Data analysis and performance studies for precision instruments in a high-requirement technical environment linked to gravitational-wave detector performance.

\textbf{Objective:} Identify, characterize and mitigate sources affecting system sensitivity using robust statistical, computational and validation-oriented methods.

\begin{itemize}
  \item \textbf{Performance validation:} Analysed sources affecting system sensitivity, performance stability and measurement quality.
  \item \textbf{Technical diagnostics:} Identified contributors to performance degradation and supported prioritisation of mitigation actions.
  \item \textbf{Analysis pipelines:} Developed Python-based pipelines for noise source identification, characterization and mitigation.
  \item \textbf{Test data analysis:} Analysed measurement and simulation outputs to identify abnormal behaviours and possible degradation sources.
  \item \textbf{First-level anomaly analysis:} Investigated unexpected behaviours, dominant contributors, weak points and possible correction or mitigation directions.
  \item \textbf{Validation reporting:} Produced reproducible and traceable technical analyses for international engineering and research stakeholders.
  \item \textbf{System impact:} Contributed to studies identifying environmental and instrumental constraints for next-generation gravitational-wave observatories.
\end{itemize}

\divider

\cvevent{\textbf{Junior R\&D Engineer - Signal Analysis, Simulation \& Software Validation}}{ARTEMIS, Observatoire de la Côte d'Azur}{2018 -- 2021}{Nice, France}

\textbf{Context:} Engineering and analysis activities for gravitational-wave mission preparation, involving signal analysis, simulation, software workflows and noise characterization.

\textbf{Objective:} Develop robust analysis methods for complex signals and support technical investigations in demanding system environments.

\begin{itemize}
  \item \textbf{Signal analysis:} Developed methods for the separation of overlapping low-SNR signals in complex datasets.
  \item \textbf{Simulation workflows:} Built simulation approaches for complex signal environments, uncertainty studies and large-scale datasets.
  \item \textbf{Software validation:} Compared simulated outputs, model behaviours and analysis results to assess consistency and robustness.
  \item \textbf{Technical investigations:} Investigated noise sources through cross-sensor correlation analyses in diagnostic contexts.
  \item \textbf{Mission performance:} Supported the identification of limiting factors affecting mission-level performance and system sensitivity.
\end{itemize}

\cvsection{Professional Training}

\cvevent{Supervised Machine Learning (Regression \& Classification)}{DeepLearning.AI - Andrew Ng}{2020}{Coursera}

\divider

\cvevent{Recherche reproductible : principes méthodologiques pour une science transparente}{FUN MOOC -- Inria}{2025}{France Université Numérique}

\divider

\cvevent{Software, Data, AI and Systems -- completed courses}{OpenClassrooms}{2024 -- 2026}{}

\small{
Python, SQL, Machine Learning, AI, Linux, JavaScript, TypeScript, Angular, Node.js, Django, REST APIs, C/C++, web development, algorithms and software engineering fundamentals.
}

\switchcolumn

\cvsection{Profile for Scalian}

\textbf{Systems validation engineer with a strong background in Python automation, IVV logic, simulation workflows, test analysis, anomaly investigation and software/hardware interface understanding}. Experience developed through major international collaborations including \textbf{LISA}, \textbf{LIGO--Virgo--KAGRA} and \textbf{Einstein Telescope}, complemented by recent industrial software experience involving operational services connected to hardware infrastructure.

For an \textbf{embedded systems validation / IVVQ / Defence} role, I bring a practical engineering mindset: developing and executing Python scripts, analysing test and simulation outputs, identifying anomalies, documenting technical evidence, supporting non-regression logic and coordinating with multidisciplinary teams.

\begin{itemize}
  \item Strong experience with complex systems under integration, validation, consistency and traceability constraints.
  \item Development and execution of Python scripts for automation, model comparison, consistency checks and validation workflows.
  \item Verification and validation of software chains, simulation workflows, model outputs and technical deliverables.
  \item Analysis of anomalies, interface behaviours, performance degradation and system-level limitations.
  \item Experience with software/hardware interfaces through operational industrial systems connected to charging infrastructure.
  \item Familiarity with technical contexts involving protocol constraints, operational data exchanges and interface reliability.
  \item Structured reporting: validation notes, technical documentation, progress summaries and technical review support.
  \item Understanding of V-cycle engineering logic: requirements, specification, integration, verification, validation and technical evidence.
  \item Experience working in Agile software environments with iterative development, task coordination and stakeholder interaction.
  \item Use of collaborative engineering tools including Confluence for technical documentation and Zenhub for Agile task tracking and coordination.
\end{itemize}

\cvsection{Fit for Embedded Systems Validation}

\begin{itemize}
  \item Relevant background for \textbf{IVVQ}, validation campaigns, non-regression logic and test result analysis.
  \item Strong practical experience with \textbf{Python automation}, scientific computing, analysis pipelines and reproducible workflows.
  \item Ability to execute scripts, inspect outputs, analyse first-level anomalies and document findings clearly.
  \item Experience producing technical documentation, validation evidence, analysis reports and review material.
  \item Comfortable with JIRA-style issue tracking, anomaly follow-up and structured technical communication.
  \item Strong fit for complex and critical environments requiring rigor, autonomy, analytical thinking and traceability.
\end{itemize}

\cvsection{Languages}

\cvskill{English}{4}
\divider
\cvskill{Spanish}{2}
\divider

\medskip

\cvsection{Education}

\cvevent{PhD in Physics}{Université Côte d'Azur}{2018 -- 2021}{Nice, France}

\divider

\cvevent{Selective Graduate Program in Fundamental Physics (Magistère)}{Université Paris-Saclay}{2016 -- 2018}{Orsay, France}

\divider

\cvevent{MSc in Astronomy and Astrophysics}{Observatoire de Paris / Université Paris-Saclay}{2017 -- 2018}{Paris, France}

\divider

\cvevent{BSc in Fundamental Physics}{Université Paris-Saclay}{2015 -- 2016}{Orsay, France}

\end{paracol}

\cvsection{Projects}

\cvevent{Co-author and full member of the LISA Consortium}{ESA/NASA Space Mission - Large-scale International Collaboration}{2018 -- now}{}

Contribution to software, analysis and validation workflows for the LISA space mission within an international engineering and scientific collaboration involving ESA, NASA and multiple research institutes.

\textbf{Role:} Software workflows, technical coordination, model integration, simulation campaigns, validation activities, complex signal-processing tasks and collaborative engineering support.

\textbf{Relevance for IVV / embedded systems validation:} Work involving complex system behaviour, mission-level performance constraints, simulation-based validation, consistency checks, non-regression logic, technical documentation, anomaly investigation and cross-functional engineering coordination.

\textbf{Program scale:} Major international space mission with an estimated budget of approximately 2 billion EUR over 10 years.

\divider

\cvevent{Co-author and full member of Einstein Telescope and LIGO--Virgo--KAGRA Collaborations}{International Collaborations}{2021 -- now}{}

Contribution to collaborative developments in data analysis, signal characterization, software methods and technical investigations within the LIGO--Virgo--KAGRA and Einstein Telescope collaborations.

\textbf{Role:} Development of analysis tools, contribution to technical activities and support to complex software workflows in high-standard collaborative environments.

\textbf{System impact:} Studies on environmental and instrumental noise sources helping identify fundamental design constraints and performance limitations for next-generation gravitational-wave observatories.
\newpage

\cvsection{Strengths}

\vspace{-.3cm}

\cvsubsection{IVVQ, validation \& test automation}

{\small
\cvtag{IVV}
\cvtag{IVVQ}
\cvtag{V-Cycle}
\cvtag{IVVQ Procedures}
\cvtag{Validation Campaigns}
\cvtag{Test Automation}
\cvtag{Python Test Scripts}
\cvtag{Documented Code}
\cvtag{Script Execution}\newline
\cvtag{Test Result Analysis}
\cvtag{Non-Regression Testing}
\cvtag{Acceptance Criteria}
\cvtag{Test Strategy}
\cvtag{Test Reports}
\cvtag{Validation Evidence}
\cvtag{User Guide Writing}
\cvtag{Progress Reporting}
}

\divider\smallskip
\vspace{-.3cm}

\cvsubsection{Systems, diagnostics \& interfaces}

{\small
\cvtag{Complex Systems}
\cvtag{Embedded Systems Context}
\cvtag{Critical Systems}
\cvtag{Defence Environment Interest}
\cvtag{System Integration}
\cvtag{Software Integration}
\newline

\cvtag{Simulation-Based Validation}

\cvtag{System-Level Analysis}
\cvtag{Anomaly Analysis}
\cvtag{First-Level Diagnostics}
\cvtag{Root Cause Analysis}
\cvtag{Hardware / Software Interfaces}
\cvtag{Interface Reliability}\newline
\cvtag{Operational Data Exchanges}
\cvtag{Protocol Constraints Awareness}
\cvtag{Data Consistency}
\cvtag{Performance Analysis}
}

\divider\smallskip
\vspace{-.3cm}

\cvsubsection{Tools, software \& coordination}

{\small
\cvtag{Python}
\cvtag{Object-Oriented Programming}
\cvtag{Linux}
\cvtag{Bash}
\cvtag{Git}
\cvtag{GitLab}
\cvtag{GitHub}
\cvtag{Docker}
\cvtag{CI/CD}
\cvtag{SQL}
\cvtag{TypeScript}
\cvtag{Angular}
\cvtag{Node.js}
\cvtag{REST API}
\cvtag{FTP}
\cvtag{JIRA}
\cvtag{Confluence}
\cvtag{Zenhub}
\cvtag{Agile Tooling}
\cvtag{Technical Coordination}
\cvtag{Stakeholder Communication}
\cvtag{Technical English}
}

\end{document}